% Working supplement. Does not replace registro_k.tex or close E108 by decree.
% Provenance and scope: INFORME_RECUPERACION.md, same folder.
\subsection{Incidence evaluation preceding the dodecaphasic record}
\label{subsec:registro-evaluacion-incidencial}

Recovering a record from its cyclic differences and
total charge is distinct from evaluating those coordinates
on the incidence ledger. The first problem is solved through the
transverse extractor of Proposition~\ref{prop:k-transversal}.
For the second, a historical presentation of the
record provides three integer counts in each window. It is useful
to make their composition explicit and determine exactly the information
their evaluation requires.

Let \(\mathscr L\) be a previously generated ledger of visits and oriented
traces. Each visit preserves the route, APP position, sheet,
instant, multiplicity, and boundary incidence. Arithmetic evaluation
produces the integer and its decomposition
\[
 n_t=r_t+9q_t,\qquad 1\leq r_t\leq9.
\]
The integer \(n_t\) is computed from the sum or product sheet; its
remainder is not received as an independent value. The distinction between
corona incidence and interior incidence remains explicit for visits
with remainder \(9\).

In window \(E_m\), the count presentation considers
\begin{align}
 A_m&=\sum_{v\in E_m}w(v)
             \mathbf1_{\{1,3,5,7\}}(r(v)),\\
 C_m&=\sum_{j\geq0}\bigl(\nu^-_{120,m}(j)-\nu^+_{120,m}(j)\bigr),\\
 V_m&=\sum_{v\in E_m}w(v)\mathbf1_{\{9\}}(r(v))\epsilon_\partial(v).
\end{align}
Here \(w(v)\) is the incidence multiplicity, and
\(\epsilon_\partial(v)=1\) on the corona, \(-1\) in the interior.
The traces counted by \(\nu^\pm\) must be enumerated through the
corresponding ledger rule. The integer sum requires finite support or an
alternative construction determining its meaning. The number of
chronological incidences and the reduced length \(120(1+3j)\) have
different types; their relation requires the reader's length unit.

Define \([x]_{10}\in\{0,\ldots,9\}\) as the Euclidean
representative and retain the quotients
\[
 x=[x]_{10}+10\lfloor x/10\rfloor.
\]
The incidence block is
\begin{equation}
 z_m=100[A_m]_{10}+10[C_m]_{10}+[V_m]_{10}.
 \label{eq:registro-bloque-incidencial}
\end{equation}
Composition yields the map
\begin{equation}
 \mathcal E_{\rm cnt}(\mathscr L)=
 \left((z_m-z_{m+3})_{m=1}^{12},
       (z_m-z_{m+4})_{m=1}^{12},\sum_{m=1}^{12}z_m\right).
\label{eq:registro-composicion-incidencial}
\end{equation}
Indices are read cyclically modulo \(12\). In this subsection,
\((D_jz)_m=z_m-z_{m+j}\), for \(j=3,4\), and
\(Q(z)=\sum_{m=1}^{12}z_m\); thus the preceding composition is
\((D_3z,D_4z,Q(z))\). This convention is retained in
\S\ref{subsec:k-transversal}, where complete integral reconstruction
is proved.
This formula produces its coordinates from the counts. It does not use
\(K\), \(U\), \(\alpha\), or the expected transverse coordinates
as arguments. Its identification with
\(\mathcal E_{108}^{90,120}\) requires composition with the effective collection and
incidences of the terminal state; it does not follow solely
from invertibility of the transverse system.

\begin{proposition}[Sufficient and necessary arithmetic information]
\label{prop:registro-36-residuos}
Let \(N,N'\in\mathbb Z^{12\times3}\), ordered by the counts
\((A,C,V)\). For the map defined by
\eqref{eq:registro-bloque-incidencial}--\eqref{eq:registro-composicion-incidencial},
\[
 \mathcal E_{\rm cnt}(N)=\mathcal E_{\rm cnt}(N')
 \quad\Longleftrightarrow\quad
 N-N'\in10\mathbb Z^{36}.
\]
In particular, the \(36\) sector residues determine exactly
the \(25\)-coordinate output. Preservation of quotients and
provenance constitutes additional information in the record.
\end{proposition}

\begin{proof}
Equality modulo \(10\) produces the same blocks and therefore
the same transverse coordinates. Conversely, if the transverse
coordinates agree, \(D_3(z-z')=D_4(z-z')=0\). Steps
\(3\) and \(4\) generate the cycle of length \(12\), so
\(z-z'\) is constant. Equality of total charges annihilates that
constant. Finally, the decimal representation of an integer between
\(0\) and \(999\) has three unique digits. Thus the
three residues of every window are recovered. The statement is an algebraic identity
for all \(N,N'\), not an inference obtained from finite tests.
\end{proof}

\subsection{Conjugation of counts and boundary terms}
\label{subsec:registro-conjugacion-incidencial}

Let \(\rho(m)=13-m\). Consider an incidence conjugation
that reverses the window order, preserves the arithmetic and
boundary classes of visits, and exchanges the trace channels.
Then
\[
 (A_m^\dagger,C_m^\dagger,V_m^\dagger)
 =(A_{\rho(m)},-C_{\rho(m)},V_{\rho(m)}).
\]
If \(c_m=[C_m]_{10}\), decimal reduction produces the identity
\begin{equation}
 z_m^\dagger=z_{\rho(m)}
       +100\mathbf1_{\{c_{\rho(m)}\ne0\}}-20c_{\rho(m)}.
 \label{eq:registro-conjugacion-decimal}
\end{equation}
Indeed, \([-C]_{10}=0\) when \([C]_{10}=0\), and
\([-C]_{10}=10-[C]_{10}\) otherwise. The formula preserves
exactly this difference. Negating the channel is not equivalent to negating
the entire decimal block.

Applying this identity to a concrete TPK involution requires
checking its action on incidences. For an occupation reader
evaluated before advancement, reversing a route
\(\gamma=(x_0,\ldots,x_n)\) satisfies
\[
 \sum_{k=0}^{n-1}f(x_{n-k})
 =\sum_{k=0}^{n-1}f(x_k)+f(x_n)-f(x_0).
\]
Endpoint terms are incorporated before reduction modulo
\(10\). On closed routes they vanish; in open windows they can
alter the residues. This correction, already present in the bidirectional
transport development, specifies what the connection
between route and counts must preserve.

\begin{proposition}[Necessity of the record's nonperiodic information]
Suppose the observable of a fixed route configuration satisfies
\(o(t+54)=o(t)\), multiplicity is preserved under that return,
and the same local reader acts in windows of \(9\) instants.
Then its counts and blocks satisfy
\[
 (A,C,V)_{m+6}=(A,C,V)_m,\qquad z_{m+6}=z_m.
\]
\end{proposition}

\begin{proof}
Translation \(t\mapsto t+54\) bijectively sends \(E_m\)
to \(E_{m+6}\), preserving each counted incidence and its multiplicity.
Equality passes to the sum and decimal reduction.
\end{proof}

Thus a presentation with \(12\) distinct blocks requires
the reader to preserve some information that does not factor through that
period-\(54\) observable: incidence selection, effective orientation,
endpoints, multiplicity, or memory. The proposition delimits an
inadmissible state reduction; it does not assert that the full dynamics has
period \(54\) or itself determine its terminal reader.
